\section{Future Directions}
\label{sec:future}

The evaluation points to a concrete programme of work, in three tracks:
near-term technical levers that follow directly from the diagnosis
(\S\ref{sec:discussion}); open design questions that need clinical partners as
much as engineering; and the prerequisites for a real field deployment. The
single highest-leverage next step is the first lever below---teaching the small
generator to use what it retrieves---which is also the experiment that would
settle why retrieval does not yet help.

\subsection{Technical levers}

Three retrieval-and-generation levers each attack a different root of the RAG
penalty, over the corpus that underlies them.

\paragraph{Generator grounding (the primary lever).} The central lever, and the
one the on-device setting most needs, is to teach the small generator to extract
specific, faithful guidance from retrieved context---to reach through retrieval
what the frontier reaches through scale, since scale is not available on the
phone (\S\ref{sec:model-size}). Concretely this is a grounding fine-tune: train
the deployed generator on (context, faithful-answer) supervision, optionally
distilled from a frontier model, so it uses on-target passages and resists
off-target ones while keeping every claim cited. The Qwen ceiling shows the
prompts and the task are already sound (\S\ref{sec:helpful-safe}); what is
missing is a small model that uses its context well. The experiment that would
settle whether RAG can be made net-positive is direct---a grounding fine-tune
evaluated under \emph{live} retrieval: if it flips RAG net-positive,
under-grounding was the cause; if a well-grounded small model still cannot beat
no-RAG, the remaining lever is capability and corpus.

\paragraph{Procedure-aware retrieval and depth.} If part of the cost is the
three-chunk window splitting a protocol, the remedy is retrieval-side, and it is
really a question of how to spend a fixed token budget bounded by latency and
the context cap. Two knobs govern how much of a procedure that budget captures:
where chunk boundaries fall (around whole protocols rather than page breaks) and
how large each chunk is, both aimed at landing a complete procedure in the
passages retrieved rather than a fragment; reranking then ensures it is the
\emph{management} section that fills the budget, not a definition or a
study-exercise (Figure~\ref{fig:case-nausea}). The budget itself is gated by
hardware: the shipped CPU default holds $k$ at three to stay within 60\,s, but
the GPU opt-in affords far more context ($k=15$ in $\sim$24\,s,
\S\ref{sec:latency}), so a faster backend---or an on-device NPU---buys procedure
coverage directly. A cheap first step is to vary depth and chunk size on the GPU
path and measure whether answer quality rises, testing the truncation hypothesis
head-on. The caution is that more retrieved text adds distraction as well as
coverage, which is why this lever and generator grounding are coupled: richer
context pays only for a generator robust enough not to drown in it.

\paragraph{Query rewriting and cross-lingual retrieval.} The cheapest lever
carries no model cost and acts where the system is most exposed: the retriever
consumes the nurse's query \emph{verbatim}. Real input is noisy---lay phrasing,
abbreviations, partial descriptions---and wording alone decides what is retrieved
(``severe nausea'' and ``hyperemesis gravidarum'' fetch different bundles for the
same patient, \S\ref{sec:retr-corpus}). A lightweight rewriting or normalisation
step that maps the query onto the corpus's vocabulary can surface the right
section before the generator runs, paired with brief user guidance on how to
describe a case. Because EmbeddingGemma is multilingual, the same front-end opens
\emph{cross-lingual} retrieval: rewriting or translating the query would let a
non-English (e.g.\ Swahili) question reach the corpus, and would let the corpus
draw on authoritative non-English references it cannot exploit today---%
high-quality German-language obstetrics texts, for one---a concrete path toward
the multilingual support a real deployment needs.

\paragraph{Corpus expansion and curation.} These three levers raise the
precision of what reaches the generator; the corpus itself remains a real
constraint. Coverage is high in aggregate---a relevant chunk exists for
$96.5\%$ of queries---but the corpus was not expanded this cycle, per-query gaps
are decisive when they fall on a query the corpus cannot answer
(\S\ref{sec:retrieval-results}), and licensing limits some high-value reference
texts. Expanding coverage and securing licensing for authoritative sources is
the slow but reliable lever, and the one whose payoff is most direct per query.

\subsection{Open design questions}

Two design questions a fielded system must answer need clinical partners as much
as engineering.

\paragraph{Conflicting guidance.} The corpus draws on many authoritative
sources, and they do not always agree. Auditing the faithfulness failures
(\S\ref{sec:generator}) surfaced 16 cases in which two retrieved passages give
conflicting advice on the same clinical question, each verified verbatim against
its source---a field reference re-doses emergency contraception after vomiting
within \emph{three} hours where WHO says \emph{two}; one WHO document starts the
active phase of labour at 4\,cm dilatation where another starts it at 5\,cm
(Appendix~\ref{app:contradictions}). The model is handed both in a single prompt,
so whichever it follows can be faulted against the other. When sources diverge,
what should the system do---prefer the local or most recent protocol, surface
both with their provenance, or flag the conflict to the clinician? Citations make
the divergence visible; choosing a behaviour, and validating it clinically, is
open for any assistant grounded in a heterogeneous guideline base.

\paragraph{Knowing when to abstain.} A deployable assistant needs a principled
way to decline, rather than guess, on a query the corpus cannot support. The
technical half is unsolved: we found no usable confidence gate today
(\S\ref{sec:retr-gate}), as neither a retrieval score nor an answer-side signal
reliably separates answerable from unanswerable queries---so a usable signal
would likely have to come from the generator (answer-side confidence or
self-consistency), which we have not built for the deployed configuration. The
harder half is clinical: \emph{when} the system should abstain, and \emph{what}
it should do instead, must be set against the nurse-midwife's responsibility
boundary---which decisions fall within their scope of practice, and which must be
escalated or referred---and defined with clinicians before any gate is tuned to
enforce it.

\subsection{Toward a field deployment}

Beyond the research levers, a real deployment has prerequisites this cycle could
not close; they bound what our evidence can claim (\S\ref{sec:limitations}) and
form the engineering-and-study track a successor would run in parallel.

\begin{itemize}\setlength{\itemsep}{2.5pt}
  \item \textbf{A field study with nurse-midwives.} Above all, the system needs
  a study with the nurse-midwives it is built for: all current evidence is
  benchmark behaviour scored by LLM judges, with no point-of-care data on real
  clinical use. This is the single most important gap.
  \item \textbf{Clinician verification.} Every quality, safety, and faithfulness
  number comes from LLM judges; the dangerous-answer determinations in particular
  await review by qualified clinicians. Independent clinical verification is a
  prerequisite for any claim of clinical correctness.
  \item \textbf{Swahili support.} The system, corpus, and evaluation are
  English-only, while the deployment setting is Swahili-speaking, and a
  safety-critical prompt cannot ship in a language it was not evaluated in
  (\S\ref{sec:system-prompt}). Validating a Swahili configuration---which the
  cross-lingual lever above complements---is a prerequisite for field use.
  \item \textbf{Validation on field hardware.} All timings come from one flagship
  SoC; the lower-cost MediaTek devices likely used in the field are untested and
  will be slower, and battery drain and sustained thermal throttling over a ward
  shift are unprofiled (\S\ref{sec:latency-pending}). The deployed G1 prompt also
  produces longer answers than the runs we timed, so measured latency is a lower
  bound.
\end{itemize}
